\section{Dynamic Wireless Topologies as Graphs}
\label{sec:background}

Fig.~\ref{fig:overview} sketches the pipeline this tutorial studies: a constellation
induces a time-varying graph, a propagation operator turns node features into a
node-level field, and the operator is the lever we vary.
\begin{figure*}[t]
\centering
\resizebox{\textwidth}{!}{%
\begin{tikzpicture}[node distance=14mm]
% Stage 1: LEO constellation (Earth + satellites on two orbits)
\node[circle, draw=steel, fill=steel!25, minimum size=11mm] (earth) at (0,0) {};
\draw[gray!55] (earth.center) ellipse (11mm and 4mm);
\draw[gray!55, rotate around={55:(earth.center)}] (earth.center) ellipse (11mm and 4mm);
\foreach \a in {15,150,255}{\fill[amber!80!black] ($(earth.center)+(\a:10.5mm)$) circle (1pt);}
\node[font=\footnotesize\bfseries, text=ink, align=center, below=7mm of earth] (l1)
  {LEO constellation\\(ephemeris-known)};

% Stage 2: time-varying graph
\node[box, minimum width=20mm, minimum height=15mm, right=18mm of earth] (g) {};
\foreach \x/\y in {-5/3,1/5,6/2,-4/-4,3/-5,6/-1}{%
  \fill[navy] ($(g.center)+(\x mm,\y mm)$) circle (1pt);}
\draw[navy!70] ($(g.center)+(-5mm,3mm)$)--($(g.center)+(1mm,5mm)$)--($(g.center)+(6mm,2mm)$);
\draw[navy!70] ($(g.center)+(-4mm,-4mm)$)--($(g.center)+(3mm,-5mm)$)--($(g.center)+(6mm,-1mm)$);
\draw[navy!70] ($(g.center)+(1mm,5mm)$)--($(g.center)+(3mm,-5mm)$);
\node[font=\footnotesize\bfseries, text=ink, align=center, below=2mm of g] (l2)
  {time-varying graph\\$G_t=(V,E_t)$};

% Stage 3: operator
\node[hot, align=center, minimum width=26mm, minimum height=15mm, right=20mm of g] (op)
  {propagation operator\\$\Prop(\Lap)$:\\GCN / heat / quantum walk};

% Stage 4: task
\node[box, align=center, minimum width=22mm, minimum height=15mm, right=18mm of op] (task)
  {node-level task\\potential / routing\\field $\hat{\mathbf{y}}$};

\draw[flow] (earth.east) -- (g.west);
\draw[flow] (g.east) -- (op.west);
\draw[flow] (op.east) -- (task.west);
\end{tikzpicture}}
\caption{System overview. A LEO constellation with known future geometry induces a
time-varying graph $G_t$; a propagation operator $\Prop(\Lap)$, classical or quantum,
turns node features into a node-level field (here a potential / routing field). This
tutorial studies the choice of $\Prop$ as the lever: the operator is the only thing that
changes between classical, diffusive, and quantum, ballistic, propagation.}
\label{fig:overview}
\end{figure*}


\subsection{The time-varying graph model}
We model a network at slot $t$ as a graph $G_t=(V,E_t)$ with $|V|=N$ nodes and a
time-varying edge set $E_t$. Let $\Adj_t\in\{0,1\}^{N\times N}$ be the adjacency
matrix (or a weighted variant carrying link delay or capacity), $\Deg_t$ the diagonal
degree matrix, and
\begin{equation}
\Lap_t=\Deg_t-\Adj_t, \qquad
\Lap_t^{\mathrm{sym}}=\Id-\Deg_t^{-1/2}\Adj_t\Deg_t^{-1/2}
\end{equation}
the combinatorial and symmetric-normalized Laplacians. Each $\Lap_t$ is symmetric
positive semidefinite, with eigendecomposition $\Lap_t=\Eig_t\,\mathrm{diag}(\lambda)\,\Eig_t^{\top}$,
$\lambda_k\ge 0$. Node attributes are collected in $\mathbf{X}\in\mathbb{R}^{N\times d}$.
Table~\ref{tab:notation} summarizes the notation used throughout.

A defining feature of LEO and other non-terrestrial
networks~\cite{kodheli2021ntn,handley2018delay} is that the future
topology is \emph{known}: orbits are deterministic over the analysis horizon, so
$\Adj_{t+1},\Adj_{t+2},\dots$ can be computed in advance from ephemeris. This is
unusual among dynamic graphs and is worth keeping in mind, although the case study
in this tutorial uses only a single snapshot per sample to keep the exposition focused
on the operator question.

\subsection{Learning tasks on $G_t$}
Many network-management problems become node-level or edge-level learning tasks once
the topology is a graph:
\begin{itemize}
\item \emph{Potential and reachability fields}: predict, for every node, a scalar
that encodes distance or reachability to a target. This underlies routing and
forwarding and is the task we use in the case study.
\item \emph{Traffic and congestion}: predict link load or a congestion price for
resource-aware forwarding.
\item \emph{Anomaly and security}: classify nodes or subgraphs as anomalous under
shifting topology.
\end{itemize}
What these share is a need to move information across the graph. How far and how fast
that information should move is exactly the property that distinguishes propagation
operators, which is why the choice of operator, not the choice of task, is the right
lever to study first.

\subsection{What makes these graphs hard}
Three features set dynamic wireless graphs apart from the static citation and molecule
graphs on which GNNs are usually taught. First, \emph{churn}: edges appear and disappear
between slots as satellites cross polar regions or seams, so an operator that overfits a
single topology generalizes poorly. Second, \emph{large diameter}: a constellation graph
is locally grid-like and globally sparse, so geodesics are long and a fixed-hop operator
is starved of range exactly where it is needed. Third, \emph{scale with structure}:
realistic shells have thousands of nodes but a regular plane-and-seam structure, which
both stresses computation and offers regularity an inductive operator can exploit. The
node-potential task we study isolates the second feature, range, because it is the one
the propagation operator most directly controls; the others motivate the roadmap in
Section~\ref{sec:roadmap}.

\begin{table}[t]
\centering
\caption{Notation used throughout the tutorial.}
\label{tab:notation}
\renewcommand{\arraystretch}{1.15}
\begin{tabular}{@{}ll@{}}
\toprule
Symbol & Meaning \\
\midrule
$G_t=(V,E_t)$ & graph at slot $t$, $N=|V|$ nodes \\
$\Adj_t,\Deg_t,\Lap_t$ & adjacency, degree, Laplacian \\
$\Eig,\lambda_k$ & Laplacian eigenvectors and eigenvalues \\
$\mathbf{X},\Feat^{(\ell)}$ & input features, layer-$\ell$ embeddings \\
$\Prop(\Lap)$ & propagation operator (the unifying object) \\
$t_\ell$ & learnable propagation time at layer $\ell$ \\
$\Wt^{(\ell)}$ & learnable layer weights \\
$\mathrm{ecc}(v),\,\mathrm{diam}(G)$ & eccentricity of $v$, graph diameter \\
\bottomrule
\end{tabular}
\end{table}
